\chapter{Vehicle Space relativ zum Eisenbahnzustand}
\label{chap:vehicle-space}

Vehicle Space führt keinen weiteren Eisenbahnzustand ein. Er fragt, wie
fahrzeugbezogene Größen relativ zum aus Teil V übernommenen qualifizierten
pose3-Frame und zur hier eingeführten Durchfahrt \(s(t)\) ausgedrückt werden.

\begin{principle}[Track--Vehicle-Unterscheidung]
Der qualifizierte Eisenbahnzustand kann ohne Fahrzeugmodell bestehen. Ein
Fahrzeugzustand benötigt diesen Bezug sowie ein identifiziertes Fahrzeugmodell
und Betriebseingaben.
\end{principle}

\begin{proposition}[Zustandshierarchie]
\label{prop:state-hierarchy}
\[
\mathrm{vehicleState}_M(t)
=
\mathcal V_M\bigl(
\mathcal S_{\mathrm{rail}}(s(t)),Q_M
\bigr),
\]
aber der Eisenbahnzustand wird nicht aus
\(\mathrm{vehicleState}_M\) abgeleitet. Modellindex und Eingaben bleiben Teil
der Ergebnisbedeutung.
\end{proposition}

\section{Relative Größen}

Das bewegte Gleis-Frame liefert longitudinale, laterale und vertikale
Bezugsrichtungen. Ein Fahrzeugmodell kann anschließend Verschiebung, Gier,
Roll, Beschleunigung, Radlast oder Umgrenzungslage eines Wagenkastens oder
Teilsystems relativ zu diesem Frame ausdrücken. Dies sind Beispiele möglicher
Modellausgaben, keine universell vorhandenen Felder; ihre Darstellung in
Vehicle Space validiert sie nicht.

\[
F_{\mathrm{vehicle}}(t)\neq
F_{\mathrm{rail}}(s(t)).
\]
Ihre relative Orientierung ist erst aussagekräftig, wenn beide Frames,
Vorzeichenkonventionen, Zeitbezug und relevante Fahrzeugfreiheitsgrade
angegeben sind.

\section{Teilsystem- und Bezugspunktverwendungen}

Radsatz, Drehgestell, Wagenkasten oder ein gewählter Schwerpunktbezugspunkt
können jeweils andere Zustandsvektoren und Modelle benötigen. Insbesondere ist
eine Schwerpunkttrajektorie nicht der gesamte Fahrzeugzustand. Sie kann nur
dann zur expliziten Optimierungsausgabe werden, wenn Fahrzeugmodell, Ziel,
Nebenbedingungen und Anwendbarkeit deklariert sind.

\begin{summary}
Vehicle Space ist eine relative Interpretationsdomäne nach qualifiziertem
Eisenbahnzustand und Durchfahrt. Er definiert weder pose3 noch Gleis-Frame oder
Alignment-Identität neu.
\end{summary}
